\documentclass[conference]{IEEEtran}
\IEEEoverridecommandlockouts

\usepackage{cite}
\usepackage{amsmath,amssymb,amsfonts}
\usepackage{algorithm}
\usepackage{algorithmic}
\usepackage{graphicx}
\usepackage{textcomp}
\usepackage{xcolor}
\usepackage{booktabs}
\usepackage{multirow}
\usepackage{url}
\usepackage{microtype}
\usepackage{balance}

\def\BibTeX{{\rm B\kern-.05em{\sc i\kern-.025em b}\kern-.08em
    T\kern-.1667em\lower.7ex\hbox{E}\kern-.125emX}}

\begin{document}

\title{MindScreen: A Tri-Modal Decision-Level Late-Fusion Framework with Deterministic Clinical Overrides for Depression Screening and Crisis Triage\\
\thanks{This research was conducted as part of the Major Capstone Project in the Department of Computer Science and Engineering, RV Institute of Technology and Management, Bengaluru, India.}
}

\author{
\IEEEauthorblockN{Savitha G\IEEEauthorrefmark{1}, Ruchita Saraf\IEEEauthorrefmark{2}, Shravya Sanikere\IEEEauthorrefmark{2}, Chandrika Lamani\IEEEauthorrefmark{2}, Apoorva K\IEEEauthorrefmark{2}}
\IEEEauthorblockA{Department of Computer Science and Engineering, RV Institute of Technology and Management, Bengaluru, India\\
Email: \IEEEauthorrefmark{1}savithag.rvitm@rvei.edu.in, \IEEEauthorrefmark{2}\{ruchitasaraf9, shravya.sanikere, chandrika.lamani, apoorva.k\}@gmail.com}
}

\maketitle

\begin{abstract}
Early identification of Major Depressive Disorder (MDD) is critical for effective therapeutic triaging, yet traditional clinical screening faces acute bottlenecks including episodic recall bias, social stigma, and severe specialist shortages. While automated screening tools have emerged, prevailing architectures operate predominantly on single modalities or deploy unconstrained neural averaging that risks diluting critical distress markers.
This paper presents \textbf{MindScreen}, an intelligent, privacy-preserving tri-modal screening framework and conversational triage companion. MindScreen introduces a decision-level late-fusion architecture synthesizing three complementary observation channels: (i) standardized clinical self-reports via the Patient Health Questionnaire-9 (PHQ-9), (ii) affective semantics derived from transformer-based natural language processing (NLP), and (iii) acoustic vocal characterization.
To resolve the vulnerability where multi-source linear averaging can conceal explicitly defined high-risk indicators, we formalize a deterministic \textbf{Clinical Safety Override (CSO)} that bypasses statistical averaging upon detection of acute clinical indicators. The CSO prevents dilution of explicitly defined high-risk indicators under three evaluated adversarial masking scenarios. Furthermore, we integrate \textbf{Saathi}, an autonomous conversational agent equipped with multi-turn context retention, affective topic classification, and crisis-resource routing with display of relevant national helpline information (Tele-MANAS and KIRAN).
Local inference executes in under 25~ms ($<$25~MB RAM), while cloud-augmented inference completes with a mean latency of 512~ms, demonstrating feasibility for resource-constrained edge and cloud healthcare deployments.
\end{abstract}

\begin{IEEEkeywords}
Multimodal late fusion, depression screening, PHQ-9, transformer NLP, acoustic biomarkers, clinical safety override, conversational triage, affective computing.
\end{IEEEkeywords}

\section{Introduction}
Major Depressive Disorder (MDD) affects over 280 million individuals globally and stands as a leading contributor to non-fatal health impairment \cite{who2023}. In lower- and middle-income nations such as India, the mental healthcare treatment gap exceeds 75\%--85\%, driven by severe shortages of licensed psychiatrists (approximately 0.75 per 100,000 population), geographic concentration of healthcare facilities in urban centers, and widespread socio-cultural stigma \cite{gururaj2016}.

The Patient Health Questionnaire-9 (PHQ-9) \cite{kroenke2001} serves as the global clinical benchmark for evaluating depressive symptom severity. However, static psychometric self-reports suffer from notable operational limitations:
\begin{itemize}
    \item \textbf{Subjective Recall Bias}: Patients frequently alter responses to conform to perceived social norms or reflect transient acute emotional states during survey completion.
    \item \textbf{Contextual Insensitivity}: Discrete Likert-scale questions cannot capture qualitative life stressors, interpersonal dynamics, or cognitive distortion patterns.
    \item \textbf{Omission of Physiological Biomarkers}: Self-reports cannot measure somatic or psychomotor correlates of depression, such as speech monotony and vocal acoustic flattening \cite{cummins2015, yang2023}.
\end{itemize}

Recent developments in affective computing have utilized transformer language models \cite{ji2022, xu2024} and vocal acoustic analysis \cite{scherer2013, cummins2015} to identify behavioral markers of depression. However, integrating heterogeneous signals into a clinical decision-support pipeline presents significant algorithmic challenges. Early-stage feature concatenation (early fusion) is highly sensitive to modality missingness (e.g., when a patient declines audio recording) and obscures the clinical explainability of individual channels. More critically, unconstrained probabilistic averaging introduces a dangerous \textit{false-negative dilution}: if a patient expresses acute suicidal intent in text or on the PHQ-9 but exhibits calm vocal acoustics, linear averaging can dilute the composite risk below the severe threshold.

To resolve these challenges, this paper presents \textbf{MindScreen}, an intelligent tri-modal screening framework combining psychometrics, transformer NLP, and acoustic analysis. The primary contributions of this work are as follows:
\begin{enumerate}
    \item \textbf{Tri-Modal Decision-Level Late-Fusion Model}: We formulate a late-fusion model ($w_T = 0.50$, $w_A = 0.30$, $w_Q = 0.20$) operating over a 4-class probability simplex, preserving modality independence and fault tolerance under missing inputs.
    \item \textbf{Deterministic Clinical Safety Override (CSO)}: We formulate an algorithmic safety boundary ensuring that explicitly defined acute indicators (PHQ-9 Item 9 $> 0$, $S \ge 20$, or crisis text tokens) bypass statistical averaging. The CSO prevents dilution of these defined indicators under adversarial conditions where one or more modalities conceal distress signals.
    \item \textbf{Safety Override Verification}: We verify the CSO on three adversarial masking scenarios (Table~\ref{tab:safety}), demonstrating that the deterministic override correctly classifies all cases that linear averaging would misclassify.
    \item \textbf{Saathi Autonomous Crisis Companion}: We develop a culturally grounded, multi-turn conversational agent with automated crisis-trigger detection and direct resource routing to Indian national helplines (Tele-MANAS and KIRAN).
    \item \textbf{Resource-Constrained Production Profiling}: We benchmark latency and memory footprint, demonstrating sub-25~ms local execution ($<$25~MB RAM) and 512~ms cloud-augmented latency.
\end{enumerate}

\section{Related Work and State-of-the-Art Taxonomy}

\subsection{Psychometric Screening Instruments}
The PHQ-9 is a 9-item validated questionnaire where items are scored from 0 to 3, yielding an aggregate score $S \in [0, 27]$. Standard clinical guidelines define five severity intervals: minimal (0--4), mild (5--9), moderate (10--14), moderately severe (15--19), and severe (20--27) \cite{kroenke2001}. In digital screening and preliminary triage workflows, categories are frequently mapped into a 4-tier decision schema (minimal, mild, moderate, severe) where scores $S \ge 15$ indicate high clinical concern \cite{kessler2003}. However, standalone questionnaire systems lack behavioral verification against somatic or vocal markers.

\subsection{Natural Language Processing in Affective Psychiatry}
Transformer architectures have achieved high benchmarks in computational psychiatry. Ji et al. \cite{ji2022} developed \textbf{MentalBERT} and \textbf{MentalRoBERTa}, demonstrating that domain-specific pre-training on mental health corpora significantly outperforms general language models on depressive symptom identification. Recently, Xu et al. \cite{xu2024} established that instruction-tuned large language models can extract nuanced cognitive distortions from patient narratives. Distilled RoBERTa models fine-tuned on emotion categorization \cite{hartmann2022} provide robust discriminative representations across valence and arousal dimensions. In production settings, pairing hosted neural transformers with deterministic lexical sentiment fallbacks is essential for guaranteeing uninterrupted clinical availability.

\subsection{Acoustic Speech Biomarkers}
Depression induces observable neuro-motor changes in speech production, including reduced vocal fold tension, restricted articulatory velocity, and flattened prosody \cite{cummins2015, scherer2013}. Benchmark studies on the Distress Analysis Interview Corpus (DAIC-WOZ) \cite{gratch2014} have shown that fundamental frequency ($F_0$) variation, Mel-Frequency Cepstral Coefficients (MFCCs), and vocal intensity serve as reliable indicators of distress. While comprehensive acoustic extraction requires significant memory, cloud-deployed screening systems benefit from calibrated vocal activity and energy representations that operate within strict microservice memory limits.

\subsection{Taxonomy of Existing Approaches}
Table~\ref{tab:taxonomy} compares MindScreen against representative state-of-the-art mental health screening systems, highlighting modality coverage, fusion paradigm, safety mechanisms, and evaluation benchmarks.

\begin{table*}[t]
\caption{Taxonomy and Comparison with State-of-the-Art Mental Health Screening Systems}
\label{tab:taxonomy}
\centering
\begin{tabular}{@{}llccclc@{}}
\toprule
\textbf{System / Framework} & \textbf{Year} & \textbf{Modalities} & \textbf{Fusion Strategy} & \textbf{Safety Override} & \textbf{Primary Focus / Limitation} & \textbf{Reported Metric} \\ \midrule
Woebot (Fitzpatrick et al.) \cite{fitzpatrick2017} & 2017 & Text Only & None (Rule/NLP) & Static Prompts & Conversational CBT; no physiological channel & PHQ-9 $\Delta = -1.72$ \\
SimSensei (DeVault et al.) \cite{devault2014} & 2014 & Video+Audio+Text & Early Feature Concatenation & No & Virtual clinical interviewer; requires lab setup & Accuracy: 74.2\% \\
Speech Survey (Cummins et al.) \cite{cummins2015} & 2015 & Audio Only & None (Acoustic Review) & No & Biomarker survey; lacks psychometric survey grounding & F1: 0.71--0.78 \\
MentalBERT (Ji et al.) \cite{ji2022} & 2022 & Text Only & None (Transformer) & No & Domain NLP pretraining; lacks safety override & Macro-F1: 0.812 \\
AVEC Benchmark (Ringeval et al.) \cite{ringeval2019} & 2019 & Audio+Video+Text & Early / Mid Fusion & No & Multimodal competition; research-only pipeline & CCC: 0.620 \\
\textbf{MindScreen (Proposed)} & \textbf{2026} & \textbf{PHQ-9+Text+Audio} & \textbf{Decision-Level Late} & \textbf{Deterministic CSO} & \textbf{Tri-modal late fusion with crisis override} & \textbf{CSO: 3/3 cases correct} \\ \bottomrule
\end{tabular}
\end{table*}

\section{System Methodology and Mathematical Formulation}

\subsection{Problem Formulation and Observation Space}
Let the input space for an assessment session be defined as a tuple $\mathcal{X} = (\mathbf{q}, T, \mathbf{a})$, where:
\begin{itemize}
    \item $\mathbf{q} = [q_1, q_2, \dots, q_9] \in \{0, 1, 2, 3\}^9$ denotes the discrete PHQ-9 survey responses,
    \item $T \in \mathcal{V}^*$ represents the free-form text journal entry, and
    \item $\mathbf{a} \in \mathcal{A}$ represents the digitized speech audio signal.
\end{itemize}

We define the discrete target space of depression severity as an ordered 4-tier set:
\begin{equation}
\mathcal{C} = \{c_0: \text{minimal}, c_1: \text{mild}, c_2: \text{moderate}, c_3: \text{severe}\}
\end{equation}
The objective of the framework is to map $\mathcal{X}$ to a calibrated probability distribution $\mathbf{p}^F \in \Delta^3$ over the 3-simplex, from which the final risk classification $\hat{y} \in \mathcal{C}$ and confidence score $c \in [0, 1]$ are derived.

\subsection{Individual Modality Processing Pipelines}

\subsubsection{Questionnaire Branch ($\mathbf{p}^Q$)}
The clinical score $S$ is calculated by summing the 9 ordinal responses:
\begin{equation}
S = \sum_{i=1}^{9} q_i, \quad S \in [0, 27]
\end{equation}
We define a piecewise mapping function $\Phi_Q: [0, 27] \to \Delta^3$ that aligns clinical cutoff intervals with the 4-class space:
\begin{equation}
\mathbf{p}^Q =
\begin{cases}
[0.80, 0.15, 0.05, 0.00]^T, & 0 \le S \le 4 \\
[0.10, 0.70, 0.15, 0.05]^T, & 5 \le S \le 9 \\
[0.00, 0.15, 0.70, 0.15]^T, & 10 \le S \le 14 \\
[0.00, 0.05, 0.15, 0.80]^T, & 15 \le S \le 27
\end{cases}
\end{equation}
In this formulation, scores $S \ge 15$ (encompassing moderately severe and severe clinical ranges) map to the severe screening tier, while preserving individual item scores for granular safety analysis.

\subsubsection{NLP Text Branch ($\mathbf{p}^T$)}
Free-text reflections $T$ are processed by an emotion classification transformer model $\mathcal{M}_{\text{NLP}}$ (DistilRoBERTa \cite{hartmann2022}), accessed via the HuggingFace Inference API. For an input sequence, the model outputs an emotion classification $e^* \in \mathcal{E}$ with softmax probability $s \in [0, 1]$:
\begin{equation}
(e^*, s) = \arg\max_{e \in \mathcal{E}} P(e \mid T; \Theta_{\text{NLP}})
\end{equation}
We define an affective translation function $\Gamma: \mathcal{E} \to \{0, 1, 2, 3\}$ based on psychological emotion-dysregulation models:
\begin{equation}
\Gamma(e^*) =
\begin{cases}
0 \ (\text{minimal}), & e^* \in \{\text{joy}, \text{neutral}, \text{surprise}\} \\
1 \ (\text{mild}), & e^* \in \{\text{fear}, \text{anger}\} \\
2 \ (\text{moderate}), & e^* \in \{\text{sadness}, \text{disgust}\}
\end{cases}
\end{equation}
Let $r = \Gamma(e^*)$. Lexical scanning checks $T$ for acute depressive tokens ($\mathcal{K}_{\text{crisis}} = \{\text{``suicide''}, \text{``kill''}, \text{``hopeless''}, \text{``worthless''}\}$). If $\exists w \in T \cap \mathcal{K}_{\text{crisis}}$, the index is adjusted: $r \leftarrow 3$.

The non-normalized score vector $\mathbf{b} \in \mathbb{R}^4$ is constructed as:
\begin{equation}
b_k = \begin{cases} s, & k = r \\ 0.10, & k \neq r \end{cases}
\end{equation}
The normalized distribution $\mathbf{p}^T$ is obtained via simplex projection:
\begin{equation}
p_k^T = \frac{b_k}{\sum_{j=0}^{3} b_j}, \quad \forall k \in \{0, 1, 2, 3\}
\end{equation}
When remote transformer inference is unreachable, an internal deterministic keyword rule engine evaluates lexical sentiment to select calibrated fallback distributions, ensuring uninterrupted clinical availability.

\subsubsection{Acoustic Processing Branch ($\mathbf{p}^A$)}
MindScreen implements a two-tier acoustic architecture addressing both research capability and cloud deployment constraints:
\begin{itemize}
    \item \textbf{Offline Research Model}: A 3-layer PyTorch MLP ($30 \to 128 \to 64 \to 4$) with Batch Normalization, ReLU, and Dropout ($p=0.3$) was designed to operate on a 30-dimensional acoustic feature vector extracted from the DAIC-WOZ corpus \cite{gratch2014}:
    \begin{equation}
    \mathbf{x}_A = [F_{0,\mu}, F_{0,\sigma}, \text{ZCR}, \text{RMS}, \boldsymbol{\mu}_{\text{MFCC}_{1\dots13}}, \boldsymbol{\sigma}_{\text{MFCC}_{1\dots13}}]^T \in \mathbb{R}^{30}
    \end{equation}
    The model checkpoint is included in the repository; large-scale prospective clinical evaluation on DAIC-WOZ is reserved for future IRB-approved trials.
    \item \textbf{Edge Production Proxy}: Under cloud-container memory limits ($<$512~MB RAM), client audio captured via HTML5 MediaRecorder (Opus/WebM) is evaluated using a speech-activity energy index $B = |\text{Base64Decode}(\mathbf{a})| / 1024$~(kB):
    \begin{equation}
    \mathbf{p}^A =
    \begin{cases}
    [0.15, 0.20, 0.55, 0.10]^T, & B < 5 \text{ kB} \ (\text{Low Activity}) \\
    [0.20, 0.55, 0.20, 0.05]^T, & 5 \le B < 30 \text{ kB} \ (\text{Restricted}) \\
    [0.65, 0.20, 0.10, 0.05]^T, & B \ge 30 \text{ kB} \ (\text{Fluent Energy}) \\
    [0.55, 0.25, 0.15, 0.05]^T, & \text{Recording Omitted}
    \end{cases}
    \end{equation}
\end{itemize}
The production energy proxy provides a clinically conservative heuristic: silence and very short clips elevate the moderate-risk prior, reducing the probability of missed moderate-to-severe presentations when acoustic data are limited.

\textbf{Clarification on which acoustic model is used in reported results:} All tri-modal fusion evaluations (Section~IV), CSO verification scenarios (Table~\ref{tab:safety}), and latency profiling (Table~\ref{tab:latency}) use the \textit{Edge Production Proxy} (Equation~9). The offline PyTorch MLP is a research artifact included in the repository for reproducibility; it is not active in the deployed cloud instance and is not used to generate any result reported in this paper.


\subsection{Decision-Level Late Fusion and Safety Constraint (CSO)}
The late-fusion module computes a convex combination of the three modality distributions:
\begin{equation}
\mathbf{u} = w_T \mathbf{p}^T + w_A \mathbf{p}^A + w_Q \mathbf{p}^Q
\end{equation}
where empirically selected clinical weights satisfy $\sum_{m} w_m = 1.0$, with $w_T = 0.50$, $w_A = 0.30$, and $w_Q = 0.20$. These weights reflect the relative clinical informativeness of each modality: free-text affective semantics are the most expressive channel, followed by vocal speech activity, with PHQ-9 providing a validated anchor. Normalizing $\mathbf{u}$ yields the fused probability vector:
\begin{equation}
p_k^F = \frac{u_k}{\sum_{j=0}^{3} u_j}, \quad \forall k \in \{0, 1, 2, 3\}
\end{equation}
The base predicted class $\hat{y}$ and confidence $c$ are determined by:
\begin{equation}
\hat{y} = \arg\max_{k \in \mathcal{C}} p_k^F, \qquad c = \max_{k \in \mathcal{C}} p_k^F
\end{equation}

\subsubsection{Clinical Safety Override (CSO)}
In clinical screening, statistical averaging must not override critical warning signs. We formalize two deterministic non-linear safety bounds:
\begin{enumerate}
    \item \textbf{Severe Score Override}: If $S \ge 20$ (the validated psychometric cutoff for severe depressive symptoms established by Kroenke et al. \cite{kroenke2001}), the classification is deterministically forced to severe. The reported confidence value is set to a rule-based priority floor---not a calibrated posterior probability---to signal the mandatory nature of the override:
    \begin{equation}
    S \ge 20 \implies \hat{y} \leftarrow \text{severe}, \quad c \leftarrow \max(0.90,\, c)
    \end{equation}
    Here, $c$ after the override denotes a minimum priority score indicating a rule-triggered outcome; it is not a statistically calibrated confidence estimate.
    \item \textbf{Acute Crisis Flag ($C_F$)}: If suicidal ideation is endorsed on PHQ-9 Question 9 ($q_9 > 0$) or acute distress keywords appear in text, emergency triage is enforced:
    \begin{equation}
    C_F = (S \ge 20) \lor (q_9 > 0) \lor (\hat{y} == \text{severe}) \lor \text{CrisisTokens}(T)
    \end{equation}
    When $q_9 > 0$ or crisis tokens are present, the system elevates triage status:
    \begin{equation}
    (q_9 > 0) \lor \text{CrisisTokens}(T) \implies \hat{y} \leftarrow \text{severe}
    \end{equation}
\end{enumerate}
Algorithm~\ref{alg:fusion} details the complete decision flow.

\begin{algorithm}[t]
\caption{Tri-Modal Decision Fusion with Clinical Safety Override}
\label{alg:fusion}
\begin{algorithmic}[1]
\REQUIRE Questionnaire answers $\mathbf{q}$, text $T$, audio $\mathbf{a}$
\ENSURE Final risk class $\hat{y} \in \mathcal{C}$, confidence $c \in [0, 1]$, crisis flag $C_F$
\STATE Calculate $S \leftarrow \sum_{i=1}^9 q_i$
\STATE Obtain $\mathbf{p}^Q \leftarrow \Phi_Q(S)$ via (3)
\STATE Obtain $\mathbf{p}^T \leftarrow \text{NLP\_Branch}(T)$ via (4)--(7)
\STATE Obtain $\mathbf{p}^A \leftarrow \text{Acoustic\_Branch}(\mathbf{a})$ via (8)--(9)
\STATE Compute weighted vector: $\mathbf{u} \leftarrow 0.50 \mathbf{p}^T + 0.30 \mathbf{p}^A + 0.20 \mathbf{p}^Q$
\STATE Normalize: $\mathbf{p}^F \leftarrow \mathbf{u} / \sum_{j} u_j$
\STATE $\hat{y} \leftarrow \arg\max_{k} p_k^F$, \quad $c \leftarrow \max_k p_k^F$
\STATE Initialize crisis flag $C_F \leftarrow \text{FALSE}$
\IF{$S \ge 20 \lor q_9 > 0 \lor \text{CheckCrisisKeywords}(T)$}
    \STATE $\hat{y} \leftarrow \text{severe}$
    \STATE $c \leftarrow \max(0.90, c)$
    \STATE $C_F \leftarrow \text{TRUE}$
    \STATE Display crisis-resource routing metadata (Tele-MANAS, KIRAN helpline numbers)
\ENDIF
\RETURN $\hat{y}, c, \mathbf{p}^F, C_F$
\end{algorithmic}
\end{algorithm}

\subsection{Saathi: Conversational Agent and Tele-MANAS Triage}
MindScreen integrates \textbf{Saathi}, an empathetic companion designed for ongoing emotional reflection. Saathi operates through a hierarchical four-tier conversational state machine:
\begin{enumerate}
    \item \textbf{Crisis Resource Router}: Detects acute self-harm ideation in dialogue and immediately displays de-escalation text alongside clickable helpline information for India's Tele-MANAS ($14416$ / $1800$-$891$-$4416$) and KIRAN ($1800$-$599$-$0019$). No automated emergency dispatch or integration with external emergency services is performed; the system surfaces information to support user-initiated contact.
    \item \textbf{Conversational Neural Core}: Employs instruction-tuned large language models (Google Gemini 1.5 Flash via REST API or Mistral-7B-Instruct via HuggingFace) parameterized with 8 turns of conversational history.
    \item \textbf{Affective Topic Classifier}: Classifies statements into six stress domains: family friction, academic pressure, insomnia, panic/anxiety, social isolation, and positive equilibrium.
    \item \textbf{Non-Pharmacological Intervention Routing}: Dynamically suggests structured somatic exercises (e.g., 4-7-8 Pranayama breathwork, 5-4-3-2-1 sensory grounding) paced every 3--4 dialogue turns to avoid interaction fatigue.
\end{enumerate}

\section{System Evaluation}

\subsection{Evaluation Approach}
MindScreen is a full-stack clinical screening prototype. System evaluation is conducted at two levels: (i) \textbf{functional correctness} of the deterministic CSO under structured adversarial scenarios, and (ii) \textbf{operational performance} of the deployment pipeline measured by latency and resource consumption. Large-scale prospective evaluation against an annotated clinical corpus requires IRB-approved data collection and is reserved for future work (see Section~\ref{sec:limitations}).

\subsection{Clinical Safety Override Verification}
The primary safety property of MindScreen is that the CSO must correctly classify all cases in which a high-risk indicator is present, regardless of what other modalities signal. To verify this property, we evaluated three adversarial masking scenarios in which linear fusion without the CSO would produce a dangerous false-negative classification (Table~\ref{tab:safety}).

\begin{table}[ht]
\caption{Safety Override Verification Under Adversarial Masking}
\label{tab:safety}
\centering
\begin{tabular}{@{}lp{2.4cm}cc@{}}
\toprule
\textbf{Test Case Scenario} & \textbf{Input State} & \textbf{Linear Fusion} & \textbf{CSO Output} \\ \midrule
TC-1: Concealed Text & $S=22$, Joy Text, Fluent Voice & Mild & \textbf{Severe ($c = 0.90$)} \\
TC-2: Item 9 Endorsed & $S=3, q_9=1$, Neutral Text & Minimal & \textbf{Severe (Crisis)} \\
TC-3: Text Suicidality & $S=2$, Crisis Word in Text & Mild & \textbf{Severe ($c = 0.90$)} \\ \bottomrule
\end{tabular}
\end{table}

In TC-1, $S = 22 \ge 20$ triggers Rule~(11), overriding linear averaging regardless of the positive text and fluent speech. In TC-2, endorsing $q_9 = 1$ triggers Rule~(13) regardless of the low total PHQ-9 score. In TC-3, detection of a crisis keyword in $T$ triggers the lexical override in the NLP branch and propagates to Rule~(13). All three cases confirm that the implemented CSO behaves deterministically as specified by Algorithm~\ref{alg:fusion}.

The CSO demonstrates correct rule execution under the three tested conditions. Broader clinical validation---covering diverse patient populations, ambiguous presentations, and spontaneous clinical interview transcripts---is explicitly identified as a requirement before operational healthcare deployment (see Section~\ref{sec:limitations}).

\subsection{Latency and Cloud Resource Footprint}
To evaluate deployment feasibility, all latency measurements were conducted on a Render.com free-tier cloud container with the following configuration:

\begin{itemize}
    \item \textbf{Hardware}: 0.5 vCPU, 512~MB RAM (shared cloud instance, x86\_64 architecture)
    \item \textbf{Software}: Ubuntu 22.04 LTS, Python 3.11, FastAPI 0.111, Uvicorn ASGI server
    \item \textbf{Concurrency}: Measurements taken under single-request sequential execution; no concurrent load was applied
    \item \textbf{Runs}: Each local microservice component was timed over 50 individual invocations following a 10-call warm-up; mean and P95 values are reported
    \item \textbf{Network latency}: Hosted API (HuggingFace Inference API, Gemini REST) latencies include round-trip network time from the cloud container to the respective API endpoint; values reflect measured wall-clock time
    \item \textbf{Memory measurement}: Peak RSS (Resident Set Size) was captured using Python's \texttt{tracemalloc} module per component invocation; values represent incremental process memory above the FastAPI baseline
    \item \textbf{Acoustic model}: The Edge Production Proxy (payload-size heuristic) was used; the offline MLP was not loaded during profiling
\end{itemize}

Table~\ref{tab:latency} reports the resulting latency and memory profile.



\begin{table}[ht]
\caption{Latency and Resource Profile Across Pipeline Microservices}
\label{tab:latency}
\centering
\begin{tabular}{@{}lccc@{}}
\toprule
\textbf{Pipeline Component} & \textbf{Mean Latency} & \textbf{P95 Latency} & \textbf{Memory Footprint} \\ \midrule
PHQ-9 Deterministic Mapping & 1.2 ms & 2.4 ms & $<1.0$ MB \\
Acoustic Energy Proxy & 8.4 ms & 14.2 ms & $<12.0$ MB \\
Local Lexical Rule Fallback & 3.6 ms & 6.1 ms & $<5.0$ MB \\
Hosted NLP (HF API) & 480.0 ms & 1,120.0 ms & External \\
Gemini 1.5 Flash (Saathi REST) & 390.0 ms & 780.0 ms & External \\ \midrule
\textbf{End-to-End Local Execution} & \textbf{13.2 ms} & \textbf{22.7 ms} & \textbf{$<$25.0 MB} \\
\textbf{End-to-End Cloud-Augmented} & \textbf{512.0 ms} & \textbf{1,150.0 ms} & \textbf{$<$85.0 MB} \\ \bottomrule
\end{tabular}
\end{table}

Local microservice execution completes in under 25~ms with an operating memory footprint under 25~MB. When augmented with remote transformer inference, the full pipeline completes with a mean latency of 512~ms (P95: 1,150~ms, depending on external cloud network latency), fully compliant with free-tier container limits (512~MB RAM).

\section{Ethical Considerations and Limitations}
\label{sec:limitations}

\subsection{Scope and Non-Diagnostic Disclaimer}
MindScreen is explicitly engineered as an \textbf{adjunctive preliminary screening and triaging tool}, not an autonomous diagnostic instrument. It does not replace comprehensive psychiatric evaluation under DSM-5 or ICD-11 criteria. Outputs are presented as risk indicators intended to facilitate timely consultation with qualified clinicians.

\subsection{Data Minimization and Privacy Architecture}
In alignment with digital healthcare data governance frameworks:
\begin{itemize}
    \item \textbf{Transient Audio Processing}: Voice recordings are processed in server memory as volatile byte streams; no raw audio files or waveforms are stored on disk or persistent databases.
    \item \textbf{Cryptographic Security}: User records utilize bcrypt password hashing ($12$ rounds) and stateless HS256-signed JSON Web Tokens (JWT).
    \item \textbf{Keyword Feature Attribution}: Rather than generating opaque predictions, the system extracts influential lexical tokens from text input, providing clinicians visibility into the risk-contributing words for each assessment.
\end{itemize}

\subsection{Limitations}
\label{subsec:limitations}
Key limitations must be clearly acknowledged:
\begin{enumerate}
    \item \textbf{Affective Proxy Mapping}: The emotion-to-severity mapping routes general affective states (sadness, fear) to depression risk tiers. This is a clinical approximation, not a validated diagnostic mapping; direct fine-tuning on labeled clinical depression corpora (e.g., DAIC-WOZ, CLPsych) would substantially improve specificity.
    \item \textbf{Acoustic Energy Proxy}: Production acoustic scoring uses speech-activity energy (Base64 decoded payload size) as a proxy for vocal biomarkers under cloud-container memory constraints. This proxy does not extract $F_0$, MFCCs, or prosodic features, and its clinical sensitivity is substantially lower than full spectrogram analysis.
    \item \textbf{No Large-Scale Clinical Validation}: The framework has not been evaluated against a labeled clinical dataset under controlled conditions. The CSO safety verification covers three deliberately constructed adversarial scenarios. Prospective clinical validation trials under IRB oversight are required before formal healthcare deployment.
    \item \textbf{Fusion Weight Justification}: The weights $w_T = 0.50$, $w_A = 0.30$, $w_Q = 0.20$ were set based on clinical reasoning about modality informativeness. Empirical optimization via cross-validation on a held-out clinical dataset is a priority for future work.
\end{enumerate}

\balance
\section{Conclusion and Future Work}
This paper presented \textbf{MindScreen}, a tri-modal decision-level late-fusion architecture for depression screening and conversational triage. By integrating validated PHQ-9 psychometric scoring, transformer-based emotional text analysis via the HuggingFace Inference API, and a speech-activity acoustic proxy, MindScreen provides a multimodal risk signal that is more informationally rich than any single modality alone. The deterministic Clinical Safety Override (CSO) prevents dilution of explicitly defined high-risk indicators (PHQ-9 Item 9 endorsement $q_9 > 0$, total score $S \ge 20$, or crisis keyword detection) by statistical averaging; this deterministic behavior was verified across three adversarial masking test cases. Integration of Saathi establishes a culturally responsive conversational companion with crisis-resource routing and display of relevant national helpline information (Tele-MANAS, KIRAN).

Future work will focus on:
\begin{enumerate}
    \item IRB-approved clinical validation trials comparing MindScreen predictions against blinded Hamilton Depression Rating Scale (HAM-D) evaluations on a prospective clinical cohort, enabling large-scale metric reporting.
    \item Replacing the acoustic energy proxy with client-side WebAudio WebAssembly (Wasm) implementations for zero-latency, on-device MFCC and $F_0$ extraction without memory constraints.
    \item Empirical optimization of fusion weights $w_T$, $w_A$, $w_Q$ via cross-validation on annotated clinical data.
    \item Longitudinal recurrent neural network modeling of daily mood trajectories to identify early warning indicators of depressive relapse.
\end{enumerate}

\begin{thebibliography}{00}
\bibitem{who2023} World Health Organization, ``Depressive disorder (depression),'' \textit{WHO Fact Sheets}, Mar. 2023. [Online]. Available: https://www.who.int/news-room/fact-sheets/detail/depression

\bibitem{gururaj2016} G. Gururaj, M. Varghese, V. Benegal, et al., ``National Mental Health Survey of India, 2015-16: Prevalence, pattern and outcomes,'' \textit{NIMHANS Publication}, no. 129, 2016.

\bibitem{kroenke2001} K. Kroenke, R. L. Spitzer, and J. B. W. Williams, ``The PHQ-9: Validity of a brief depression severity measure,'' \textit{Journal of General Internal Medicine}, vol. 16, no. 9, pp. 606--613, Sep. 2001.

\bibitem{cummins2015} N. Cummins, S. Scherer, J. Krajewski, S. Schnieder, J. Epps, and T. F. Quatieri, ``A review of depression and suicide risk assessment using speech analysis,'' \textit{Speech Communication}, vol. 71, pp. 10--49, Jul. 2015.

\bibitem{ji2022} S. Ji, T. Zhang, L. Ansari, J. Fu, P. Tiwari, and E. Cambria, ``MentalBERT: Publicly available pretrained language models for mental healthcare,'' in \textit{Proc. 13th Language Resources and Evaluation Conference (LREC)}, 2022, pp. 7184--7190.

\bibitem{xu2024} X. Xu, B. Zou, Y. Ding, et al., ``Mental-LLM: Leveraging large language models for mental health prediction,'' in \textit{Proc. 62nd Annual Meeting of the Association for Computational Linguistics (ACL)}, 2024, pp. 5120--5135.

\bibitem{yang2023} L. Yang, D. Jiang, and E. Cambria, ``A survey on multimodal depression detection,'' \textit{IEEE Transactions on Affective Computing}, vol. 14, no. 4, pp. 3125--3144, 2023.

\bibitem{scherer2013} S. Scherer, G. Stratou, M. Mahmoud, et al., ``Automatic behavior descriptors for psychological disorder analysis,'' in \textit{Proc. 10th IEEE International Conference on Automatic Face and Gesture Recognition (FG)}, 2013, pp. 1--8.

\bibitem{hartmann2022} J. Hartmann, ``Emotion English DistilRoBERTa-base,'' \textit{Hugging Face Model Hub}, 2022. [Online]. Available: https://huggingface.co/j-hartmann/emotion-english-distilroberta-base

\bibitem{gratch2014} J. Gratch, R. Artstein, G. Lucas, et al., ``The Distress Analysis Interview Corpus of human and computer interviews,'' in \textit{Proc. 9th International Conference on Language Resources and Evaluation (LREC)}, 2014, pp. 3123--3128.

\bibitem{kessler2003} R. C. Kessler, P. R. Barker, L. J. Colpe, et al., ``Screening for serious mental illness in the general population,'' \textit{Archives of General Psychiatry}, vol. 60, no. 2, pp. 184--189, Feb. 2003.

\bibitem{fitzpatrick2017} K. K. Fitzpatrick, A. Darcy, and M. Vierhile, ``Delivering cognitive behavior therapy to young adults with symptoms of depression and anxiety using a fully automated conversational agent (Woebot): A randomized controlled trial,'' \textit{JMIR Mental Health}, vol. 4, no. 2, p. e19, Jun. 2017.

\bibitem{devault2014} T. DeVault, R. Artstein, G. Benn, et al., ``SimSensei Kiosk: A virtual human interviewer for healthcare decision support,'' in \textit{Proc. 2014 International Conference on Autonomous Agents and Multi-agent Systems (AAMAS)}, 2014, pp. 1061--1068.

\bibitem{ringeval2019} F. Ringeval, B. Schuller, M. Valstar, et al., ``AVEC 2019 workshop and challenge: state-of-mind, detecting depression with AI, and cross-cultural affect recognition,'' in \textit{Proc. 9th International Audio/Visual Emotion Challenge and Workshop}, 2019, pp. 3--12.

\bibitem{telemanas2022} Ministry of Health and Family Welfare, Government of India, ``Tele Mental Health Assistance and Networking Across States (Tele-MANAS),'' \textit{National Health Mission}, 2022. [Online]. Available: https://telemanas.mohfw.gov.in
\end{thebibliography}

\end{document}
